\clearpage
\section{The affine-vector commutant and its Fourier-polarized form}
\label{sec:f1-limit-mirabolic}

Adding a vector to a pair of flags produces a finite arithmetic observable
algebra that remembers more than the ordinary flag Hecke algebra.  Its extra
generator is a controlled translation projector.  Fourier transformation
turns these translation constraints into controlled phase constraints, but
the flag register must simultaneously be dualized and reversed.  We fix all
indices and signs below, including the square of the Fourier bridge.

\subsection{Affine invariant kernels}

\begin{definition}[Arithmetic affine-vector context algebra]
Let $k=\mathbb F_Q$, let $L$ be a finite-dimensional $k$-vector space,
$G=GL(L)$, and $X=\operatorname{Fl}(L)$ the set of complete flags.  The
affine group
\[
 \operatorname{Aff}(L)=G\ltimes L,
 \qquad (g,a)(h,b)=(gh,a+gb),
\]
acts on $Y_L=X\times L$ by
\[
 (g,a)(F,x)=(gF,a+gx).
\]
Define
\[
 R_X(L)=\operatorname{End}_{\operatorname{Aff}(L)}\mathbb C[Y_L]
\]
with Hilbert adjoint and normalized operator trace
$\tau_Q=\operatorname{Tr}/(|X||L|)$.  The invariant-kernel convention is
\[
 K_f((F,x),(F',x'))=f(F,F',x'-x).
\]
\end{definition}
\begin{scope}
The field cardinality is denoted $Q$.  The symbol $q$ below is reserved for a
real deformation parameter.  The affine group symbol does not reuse the
notation for a Weyl subalgebra.
\end{scope}
\provenance{D1241}{theory/sidequests/f1-limit/mirabolic.md}{theory/checks/f1\_limit\_mirabolic\_check.py (B1)}{theory/verdicts/f1-limit-adjudication.md}

\begin{theorem}[Affine convolution as a physical commutant]
{\normalfont\statusProved}\quad
For every finite field $\mathbb F_Q$ and finite-dimensional $L$, convolution
on $GL(L)$-orbits of triples $(F,F',v)$ is *- and trace-preservingly the
commutant $R_X(L)$.  In particular its normalized operator trace is faithful
and positive.
\end{theorem}
\begin{scope}
This is an arithmetic finite-Hilbert-space theorem.  It does not infer
positivity of a generic deformation from a polynomial presentation.
\end{scope}
\provenance{F1-MIR-AFF}{theory/sidequests/f1-limit/mirabolic.md}{theory/checks/f1\_limit\_mirabolic\_check.py (B1, n=2 and Q=2,3)}{theory/verdicts/f1-limit-adjudication.md}
\begin{proof}
If $B\leq GL(L)$ stabilizes a reference complete flag, then
$\operatorname{Aff}(L)/(B,0)\cong X\times L$ by $(g,a)B\mapsto(gB,a)$.
An operator commutes with the affine action exactly when its matrix kernel is
constant on diagonal affine orbits.  Translation by $-x$ reduces a pair to
$((F,0),(F',x'-x))$, leaving diagonal $GL(L)$-invariance.  Matrix
multiplication becomes
\[
 (f*g)(F,F',v)=\sum_{H\in X}\sum_{u\in L}
                   f(F,H,u)g(H,F',v-u),
\]
which is the stated convolution.  Taking adjoints gives
$f^*(F,F',v)=\overline{f(F',F,-v)}$.  Normalized matrix trace restricts to a
faithful positive trace on every finite-dimensional *-subalgebra.
\end{proof}

\subsection{The controlled translation generator}

For $u\in L$ and an additive character $\chi$ of $L$, write
$X(u)e_x=e_{x+u}$ and $Z(\chi)e_x=\chi(x)e_x$.  For a subspace $U\leq L$,
put
\[
 P_U^X=|U|^{-1}\sum_{u\in U}X(u),\qquad
 P_U^Z=|U^{\mathrm{ann}}|^{-1}
       \sum_{\chi\in U^{\mathrm{ann}}}Z(\chi).
\]

\begin{definition}[Controlled translation and phase constraints]
For a complete flag $F=(U_i(F))_{i=0}^n$, with $U_0=0$ and $U_n=L$, define
on $\mathbb C[\operatorname{Fl}(L)]\otimes\mathbb C[L]$
\[
 C_i^X=\sum_F|F\rangle\!\langle F|\otimes P_{U_i(F)}^X,
 \qquad
 C_i^Z=\sum_F|F\rangle\!\langle F|\otimes P_{U_i(F)}^Z,
 \qquad 0\leq i\leq n.
\]
Thus $C_0^Z=I\otimes|0\rangle\!\langle0|$ and $C_n^Z=I$.  Let
$\operatorname{Ctx}_X(L)$ be the *-algebra generated by all relative-position
flag operators $A_w\otimes I$ and all $C_i^X$.  For $n\geq1$, its marked
projection and unnormalized vector generator are
\[
 e=C_1^X,\qquad T_0=Qe-I.
\]
The same definitions apply on the dual vector space.
\end{definition}
\begin{scope}
The endpoints $i=0,n$ are included.  In particular the index-zero phase
constraint is non-scalar in rank one.  A named additive character is used
when linear duals are identified with additive characters.
\end{scope}
\provenance{D1142, D1242}{theory/sidequests/f1-limit/mirabolic.md}{theory/checks/f1\_limit\_mirabolic\_check.py (B2); theory/checks/f1\_limit\_bridge\_check.py (B3)}{theory/verdicts/f1-limit-adjudication.md}

\begin{theorem}[Generation by flags and one controlled line]
{\normalfont\statusProved}\quad
For $L$ of dimension $n\geq1$ over $\mathbb F_Q$, the embedded flag Hecke
algebra acts as $A_w\otimes I$, the marked element
$(T_0+I)/Q$ is exactly $C_1^X$, and these operators generate the full affine
commutant $R_X(L)$.  Moreover
\[
 e^2=e=e^*,\qquad \tau_Q(e)=Q^{-1},\qquad
 T_0^2=(Q-2)T_0+(Q-1)I.
\]
\end{theorem}
\begin{scope}
The additional constraints $C_i^X$ are useful tests but are redundant as
algebra generators.  The theorem fixes the orbit basis, adjoint, trace and
Hecke inclusion, rather than asserting only an abstract algebra isomorphism.
\end{scope}
\provenance{F1-MIR-GEN}{theory/sidequests/f1-limit/mirabolic.md}{theory/checks/f1\_limit\_mirabolic\_check.py (B2, n=2 and Q=2,3)}{theory/verdicts/f1-limit-adjudication.md}
\begin{proof}
The zero-vector orbit kernels act as $A_w\otimes I$.  The kernel of
$T_0+I$ is one precisely when the two flags agree and the vector difference
lies in the controlled line $U_1(F)$.  Dividing by $Q$ averages translations
along that line, hence gives $C_1^X$.  A finite-group average is a
self-adjoint projection; its rank in each flag block is $|L|/Q$, which gives
the trace.  The quadratic formula follows from $T_0=Qe-I$.  Rosso's
generation theorem then gives the reverse inclusion from these concrete
operators to the full invariant-kernel algebra.
\end{proof}

\subsection{Fourier duality with exact indices}

\begin{definition}[Fourier dual flag reversal]
Fix a named nontrivial additive character $\psi:k\to U(1)$.  Define
\[
 \mathcal F_{L,\psi}e_x
   =|L|^{-1/2}\sum_{\lambda\in L^\vee}
       \psi\bigl(-\lambda(x)\bigr)e_\lambda.
\]
If $F=(U_i)_{i=0}^n$, its reversed dual flag is
\[
 D_L(F)_j=U_{n-j}^{\perp},\qquad 0\leq j\leq n.
\]
Let $D_L$ also denote the induced flag-basis unitary and put
$W_{L,\psi}=D_L\otimes\mathcal F_{L,\psi}$.  Define
\[
 R_Z(L^\vee)=W_{L,\psi}R_X(L)W_{L,\psi}^*.
\]
The phase-polarized groupoid has finite-dimensional $k$-spaces as objects,
linear isomorphisms as arrows, and retains the pair $(R_X(L),R_Z(L^\vee))$,
their Hecke inclusions, all controlled constraints and the named $W_{L,\psi}$.
A same-register realization additionally names a linear self-duality
$\sigma:L\to L^\vee$; morphisms act on flags, vectors and duals by their
induced permutation unitaries.
\end{definition}
\begin{scope}
The negative Fourier kernel, the phase $\psi$, dual-flag reversal and any
same-register self-duality are named choices.  Forgetting them retains only
the corresponding conjugacy class of algebra comparisons.
\end{scope}
\provenance{D1243, D1244}{theory/sidequests/f1-limit/mirabolic.md}{theory/checks/f1\_limit\_bridge\_check.py (B3)}{theory/verdicts/f1-limit-adjudication.md}

\begin{theorem}[Fourier exchange of controlled constraints]
{\normalfont\statusProved}\quad
For every finite field $k=\mathbb F_Q$, every $n$-dimensional $L$, every
named nontrivial $\psi:k\to U(1)$ and every $0\leq i\leq n$,
\begin{align*}
 W_{L,\psi}C_i^XW_{L,\psi}^*&=C_{n-i}^Z,\\
 W_{L,\psi}(A_w\otimes I)W_{L,\psi}^*
   &=A_{w_0ww_0}\otimes I.
\end{align*}
The operators on the right generate exactly $R_Z(L^\vee)$, including
$C_0^Z=I\otimes|0\rangle\langle0|$.  Under canonical double duality,
\[
 W_{L^\vee,\psi}W_{L,\psi}e_{(F,x)}=e_{(F,-x)}.
\]
The construction is natural for linear isomorphisms.
\end{theorem}
\begin{scope}
Only the vector Fourier factor is tensor-preserving under direct sums.  The
full flag--Fourier bridge is not monoidal on complete-flag registers and does
not identify either commutant with the full Weyl matrix algebra.
\end{scope}
\provenance{F1-MIR-FOURIER}{theory/sidequests/f1-limit/mirabolic.md}{theory/checks/f1\_limit\_bridge\_check.py (B3, Fourier constraints, reversal, square and ranks)}{theory/verdicts/f1-limit-adjudication.md}
\begin{proof}
Character orthogonality makes $\mathcal F_{L,\psi}$ unitary and gives
\[
 \mathcal F_{L,\psi}P_U^X\mathcal F_{L,\psi}^*=P_{U^\perp}^Z.
\]
Annihilators reverse a full flag, so the index $i$ becomes $n-i$.  The
intersection-rank matrix of a flag pair transforms by
$r'(i,j)=i+j-n+r(n-i,n-j)$, which is the rank matrix of $w_0ww_0$.
This proves both conjugation formulas and the generation statement.
Applying the negative Fourier kernel twice yields
\[
 |L|^{-1}\sum_{y\in L}\sum_{\lambda\in L^\vee}
       \psi\bigl(-\lambda(x+y)\bigr)e_y=e_{-x};
\]
dual reversal squares to the original flag.  In characteristic two this
vector negation is the identity.
\end{proof}

\begin{definition}[Vector-polarization bridge package]
In rank $n$, a bridge package retains the affine commutant and orbit data,
the controlled translation and phase families for every $0\leq i\leq n$,
the based Hecke inclusion, conjugate-linear orbit star, reference trace,
algebraic and trace-supported endpoints, the named phase and Fourier unitary,
and the regular one-sided operational boundary.  When composition is needed,
it additionally names an arithmetic symplectic or affine decomposable-shuffle
correspondence with its ternary shuffle coherence.

Only $\mathcal F_{L\oplus M}=\mathcal F_L\otimes\mathcal F_M$ is a literal
tensor identity.  On the decomposable flag--vector subspace the full bridge
uses
\[
 W_{L\oplus M}J_{L,M}
 =J_{L^\vee,M^\vee}
   \bigl((W_L\otimes W_M)\otimes\operatorname{Rev}\bigr),
\]
where $\operatorname{Rev}$ reverses the shuffle word.  An isomorphism of
packages preserves all named UCP maps, the regular label class and quotient
boundary functor, and the exact square
$W_{L^\vee,\psi}W_{L,\psi}(F,x)=(F,-x)$.
\end{definition}
\begin{scope}
An abstract algebra isomorphism is not an isomorphism of bridge packages.
The decomposable flag subspace is generally proper; even adjoining shuffles
does not tensor-factor the full complete-flag Hilbert space.
\end{scope}
\provenance{D1259}{theory/sidequests/f1-limit/mirabolic.md; theory/sidequests/f1-limit/type-c.md}{B1--B7 in the two bridge checkers, at their stated finite scopes}{theory/verdicts/f1-limit-adjudication.md}
